
%forthe precision \( \gamma = \frac{1}{\sigma^2} \):

State the classical normal linear regression model. Which parameters are the
objective of (Bayesian) statistical inference?

The two-parameter model described by the conditional probability mass function
\[
f_{y | \beta, \sigma^2, X}(y | \beta, \sigma^2, X) = (2\pi\sigma^2)^{-n/2} \exp \left( -\frac{1}{2\sigma^2} \sum_{i=1}^{n} (y_i - \beta^T X_i)^2 \right)
\]
\[
\propto \exp \left( -\frac{1}{2\sigma^2} [y^T y - 2\beta^T X^T y + \beta^T X^T X \beta] \right)
\]
i.e., \(y | \{\beta = \beta, \sigma^2 = \varsigma^2, X = X\} \sim N(X\beta, \varsigma^2 I)\) is called the Bayesian normal linear regression model.


Which marginal prior distributions allow – assuming a product prior – a tractable
computation of the full conditional distributions?

Since \( \mathbf{y} \) is normal, the same could be assumed about \( \mathbb{\beta} \), say \( \mathcal{N}(\beta_0, \Sigma_0) \), which indeed
gives a multivariate normal as posterior.

\(f_{\beta | y, X, \sigma^2}(\beta | y, X, \sigma^2) \propto f_{y | X, \beta, \sigma^2}(y | X, \beta, \sigma^2) \cdot f_{\beta}(\beta)\)
\[
\propto \exp \left( \beta^T [\Sigma_0^{-1} \beta_0 + X^T y / \varsigma^2] - \frac{1}{2} \beta^T [\Sigma_0^{-1} + X^T X / \varsigma^2] \beta \right),
\]
which is proportional to a multivariate normal density \(N(\beta_n, \Sigma_n)\), where
\[
\Sigma_n = (\Sigma_0^{-1} + X^T X / \varsigma^2)^{-1} \quad \beta_n = \Sigma_n (\Sigma_0^{-1} \beta_0 + X^T y / \varsigma^2).
\]

If the elements of the prior precision matrix \(\Sigma_0^{-1}\) are small in magnitude, the conditional expectation \(\beta_n\) is approximately 
equal to the least squares estimator \(LSE = (X^T X)^{-1} X^T y\) (the estimator from the data).
If, on the other hand, the sampling precision is very small (i.e., \(\varsigma^2\) is very large), the expectation is approximately 
equal to the prior expectation \(\beta_0\).

Analogous to before we choose for \( \sigma^2 \) an inverse-gamma distribution, which yields as posterior distribution 

\( \{ \sigma^2 | \mathbf{y} = y, \mathbf{X} = X, \mathbf{\beta} = \beta \} \sim \text{IG}(\frac{\nu_0 + n}{2}, \frac{\nu_0\simga_0^2 + \text{LS}(\beta)}{2} \)

where \(\text{LS}(\beta)\) denotes the sum of squared deviations.

Together they can be estimated with a Gibbs sampler in the same fashion as we did for the normal model.

Why, in particular, are often priors important that incorporate significant uncer-
tainty? How can this be implemented? Discuss why MCMC is typically necessary to compute the posterior. 

Since priors are hard to compute (the number of prior correlation parameters grows with \( \binom{p}{2} \), \( p \) the number of covariates).
If the prior distribution is not going to represent real prior information about the parameters, then it
should be as minimally informative as possible (→ weakly or non-informative priors).
The resulting posterior distribution would then represent the posterior information of someone who began with
little knowledge of the population being studied.
Example:
⋆ Idea: Increase variances of normal prior for \( \beta \) to incorporate uncertainty of prior knowledge.
→ Limiting case: \(f_\beta \propto \text{const.}\)
⋆ Such a prior does not constitute a probability density! It is called improper or diffuse.
⋆ Still, diffuse priors are admissible as long as the resulting posterior distribution is a proper distribution. In this special case,
the posterior is proper, but does not possess a tractable parametric form.
→ The analysis can be based on MCMC methods.


Explain the idea of Bayesian model selection. How does the Bayesian formula for the indicator vector look like? 

The question continues with:
Describe the steps in the Gibbs sampler.

To reduce the number of covariates for which priors have to be computed, we can enlarge the model as follows:
Write regression coefficients as \(\beta_j = z_j \cdot b_j\), where \(z_j \in \{0, 1\}\) and \(b_j \in \mathbb{R}\), \(j = 1, ..., p\).
The regression equation becomes
\[
y_i = z_1 b_1 x_{i,1} + ... + z_p b_p x_{i,p} + \varepsilon_i, \quad i = 1, ..., n,
\]
i.e., the \(z_j\)'s indicate which regression coefficients are non-zero.
Each value of \(z = (z_1, ..., z_p)\) corresponds to a choice of covariates to include (which is the model selection)

Given a prior distribution \(f_z\) over models, we can compute a posterior probability for each regression model:
\(
f_{z | y, X}(z | y, X) = \frac{f_z \cdot f_{y | X, z}}{f_{y | X}(y | X)} = \frac{f_z \cdot f_{y | X, z}}{\sum_{\tilde{z}} f_z(\tilde{z}) \cdot f_{y | X, \tilde{z}}}
\)
if we evaluate \(f_{y | X, z}\) for each possible set of covariates \(\tilde{z}\).

The corresponding Gibbs sampler (because the sum in the denominator consists of \( 2^p \) terms) is

Input: Initial values \(\{z^{(0)}, \sigma^{2,(0)}, \beta^{(0)}\}\); Set \(s = 0\).

Algorithm:
Step 1: Set \(z = z^{(s)}\);
Step 2: For \(j \in \{1, ..., p\}\) in random order, replace \(z_j\) with a sample from \(f_{z_j | z_{-j}, y, X}\) 
determined using marginal probabilities \(f_{y | X, z}\)
Step 3: Set \(z^{(s+1)} = z\);
Step 4: Sample \(\sigma^{2,(s+1)} \sim f_{\sigma^2 | z^{(s+1)}, y, X}\); [inverse gamma]
Step 5: Sample \(\beta^{(s+1)} \sim f_{\beta | z^{(s+1)}, \sigma^{2,(s+1)}, y, X}\); [multivariate normal]
Step 6: Set \(s = s + 1\) and return to Step 1;
Output: Dependent samples \(\{z^{(s)}, \sigma^{2,(s)}, \beta^{(s)}\}, s = 1, 2, ...\)
